Un taller fabrica \textbf{mesas} ($x_1$) y \textbf{sillas} ($x_2$). Cada mesa deja
un beneficio de 5 u.m. y cada silla de 4 u.m. La carpintería dispone de 18 horas
y el acabado de 12 horas; una mesa consume 3 horas de carpintería y 1 de acabado,
y una silla 2 y 2 respectivamente. Por contrato deben fabricarse al menos 2 mesas.
El modelo que maximiza el beneficio es el siguiente:

\begin{equation}
\begin{split}
\label{eq:demo_muebles}
\mbox{max. } z = 5x_1 + 4x_2\\
s.a.:\\
3x_1 + 2x_2 \leq 18\\
1x_1 + 2x_2 \leq 12\\
1x_1 \geq 2\\
x_1,\,\,x_2\geq 0\\
\end{split}
\end{equation}

\begin{enumerate}
    \item Formular el problema de la primera fase del método de las dos fases.
    \item Resolver el problema mediante el método de la matriz completa.
    \item Indicar el plan de producción óptimo y el beneficio alcanzado.
\end{enumerate}
